\documentclass[12pt]{article}
\usepackage{kyradefaults}
\addbibresource{references.bib}

\title{\textbf{All Aboard? Causal Evidence on the Labor-Market Effects of New Rail Stations}}
\author{Kyra Sadovi%
  \thanks{E-Mail: \texttt{ksadovi@uchicago.edu}}}
\affil{Harris School of Public Policy, University of Chicago}
\date{}

\begin{document}
\maketitle
\thispagestyle{empty}

Public transit is among the most consequential infrastructural interventions a city can make, yet credible causal evidence on its effects on workers' economic outcomes remains limited. This paper asks: does gaining physical access to a new rail station improve the labor-market outcomes of nearby residents? I focus on two outcomes---worker flows and income---and argue that a new station can affect both through three mechanisms: enabling workers to more reliably retain existing jobs, expanding their effective job-search radius to include higher-paying opportunities across the metro area, and widening their referral networks by connecting them to a geographically broader set of social ties. Together, these mechanisms suggest that transit access is not merely a convenience amenity but a potential driver of upward economic mobility---one distributed very unequally across urban space.

\section{Identification Strategy}

Identifying the causal effect of transit access is nontrivial. Three identification challenges complicate standard approaches. First, station locations are not randomly chosen: municipalities tend to build in dense, well-resourced, or politically influential neighborhoods, or in areas already experiencing labor-market growth. Second, residents sort endogenously into neighborhoods near planned stations, meaning that those who eventually live near new stations may differ systematically from those who do not. Third, transit projects have long planning horizons---often a decade or more---creating substantial anticipatory effects in which residents and firms adjust behavior long before a station opens.

I address all three challenges by exploiting quasi-random variation in the \textit{timing} of station openings generated by construction delays. My identification assumption is that, conditional on a station having been selected for a particular neighborhood, the length of the delay is as good as random---driven by engineering complications, supply chain disruptions, permitting bottlenecks, and other factors unrelated to local labor-market trends. By working only with neighborhoods already selected to receive a station and varying only \textit{when} that station opens, I implicitly identify the average treatment effect on the treated: the effect of transit access on residents in neighborhoods that were chosen for a station. This sidesteps the station-placement endogeneity problem without requiring me to model why cities choose particular locations.

To study how effects evolve over time and to account for staggered adoption across projects, I embed this design into a stacked difference-in-differences framework. For each adoption cohort---defined by the realized opening date of a station---I form a sub-experiment comparing nearby tracts to tracts that remain untreated throughout the corresponding event window. Stacking these sub-experiments yields a balanced event-study design in which changes in estimated effects across event time reflect true dynamics rather than shifts in cohort composition. Critically, using planned opening dates from Environmental Impact Statements as my baseline time index allows me to test for anticipatory effects: I can examine whether residents begin adjusting labor-market behavior before a station opens, and how excess delays dampen or distort those adjustments.

\section{Data}

A central contribution of this paper is the construction of a novel dataset of U.S.\ rail transit projects since 2000, linking planned opening dates, realized opening dates, and station geographies for every major rail project in my sample. Planned opening dates are derived from Draft Environmental Impact Statements (DEIS) archived in the EPA's EIS database---specifically, the first projected completion date appearing in each project's DEIS. This ensures that my ``planned'' dates reflect genuinely ex-ante projections rather than revised estimates produced after delays were already underway. My current sample covers projects in New York City, Washington D.C., San Francisco, Los Angeles, Seattle, Boston, Miami, and Honolulu, supplemented by data from the Transit Costs Project at NYU \parencite{levy_2025_9wnjp-kez15}.

Labor-market outcomes are measured using the LEHD Origin-Destination Employment Statistics (LODES), a Census Bureau dataset recording job flows between residential and workplace Census tracts. I use these data to construct measures of worker outflows---tracking how many workers living in a given tract are employed in each destination tract---and follow how these flows change as new stations open nearby.

Transit exposure is operationalized using walking-time isochrones computed via the \texttt{r5r} package \parencite{r5r}, which builds a multimodal routing graph from OpenStreetMap pedestrian networks and GTFS feeds. For each Census tract centroid, I compute walking times to each station under the pre-period street network, assigning tracts to proximity bands of 0--5, 5--10, and 10--15 minutes. This replaces cruder Euclidean distance measures and accounts for real-world barriers---rivers, highways, block structures---that shape whether residents can realistically reach a station on foot.

\section{Contribution}

This paper provides new causal evidence on how public transit shapes labor-market mobility in American cities, with direct implications for policy debates about transit investment, urban inequality, and the geographic concentration of economic opportunity. By constructing a novel EIS-based delay dataset and combining it with a stacked event-study design, it offers a replicable framework for evaluating large infrastructure interventions in settings where project selection is non-random but completion timing is not.

\newpage
\printbibliography

\end{document}